\chapter{Parameters to set and lock}
\label{chap:parameters}

Chapter \ref{chap:criteria} says what will be judged. This chapter says what
has to be set, in what order, and why each setting is what it is. The numeric
thresholds are not repeated here; they are in Tables \ref{tab:crit-calib} to
\ref{tab:crit-record}.

Each entry is marked \textbf{Physics} if the Sun or the sensor fixes it and
the project cannot negotiate, or \textbf{Decision} if it is a project choice
that could be revised with an argument.

\section{Illumination}

\subsection*{L1. One direction of light. \textnormal{Physics}}

In orbit the tether is lit by a single distant source and there is no
atmosphere to fill the shadows: the lit side is at full sun and the shadowed
side is nearly black. Several fixtures are permitted only if their axes are
parallel. Converging beams are forbidden, and this is the most important
prohibition in the chapter: two converging sources put two shadows and two
specular highlight lines on a target that is one or two pixels wide, which
presents the matching stage with a doubled thread, a failure mode it is not
designed to handle. Also forbidden are mixing fixture types or colour
temperatures along a row, and placing a fixture at one end shining along the
tether, which combines grazing incidence with inverse-square falloff along
its length.

\subsection*{L2. Daylight colour. \textnormal{Decision}}

Every space vision laboratory in the reviewed set that declares a colour
temperature falls in the daylight range, and the DLR EPOS sun simulator is a
6000~K HMI with CRI 90 \cite{dlr2025eposportfolio, benninghoff2017epos}. The
Sun's effective temperature is 5772~K \cite{prsa2016iau}. This rules out
3200~K tungsten without further discussion. TLCI is requested as a datum
because CRI was not designed for cameras \cite{ebu_tech3355}, but no TLCI
threshold is set, because no reviewed source gives one.

\subsection*{L3. Parallel beam. \textnormal{Decision, anchored to physics}}

Verified with the penumbra test of Section \ref{sec:penumbra}. There is a
second, independent reason for collimation beyond matching the Sun's shadow
behaviour: without it, the inverse-square falloff makes the uniformity
requirement L4 unachievable at any distance that fits in the room. A point
source would need a throw of roughly twenty times the object depth to reach
class B uniformity, which is about 10~m for a 0.5~m deep structure and 20~m
for a 1~m one. Few rooms offer that.

\subsection*{L4. Uniformity across the working volume, depth included.
\textnormal{Decision, from IEC class B}}

Measured on a lux grid that spans width, height and depth. The depth axis
fails first and is the one most often left out.

\subsection*{L5. Enough light at the tether. \textnormal{Physics, from the
datasheet}}

Derived in the dark-room deliverable from the IMX477 datasheet. It is a
floor. Four assumptions enter the derivation, all of which push the true
requirement upward: the datasheet's 3200~K test spectrum differs from the
5600 to 6000~K in use; the treatment of sensor gain; an on-axis transmittance
of unity, ignoring vignetting; and a white card of reflectance 0.9 filling
many pixels, where the real tether is darker and narrower.

\subsection*{L6. No flicker. \textnormal{Physics}}

For an HMI, use an electronic ballast in its high-frequency mode, near
1000~Hz. Note specifically that the ARRI EB 12/18 manual states its
``50/60~Hz, low noise'' mode is \emph{not} flicker free despite the ballast
being an electronic one, and that its 75~Hz ``flicker free'' mode is also
unsuitable here because the resulting period exceeds the exposure ceiling set
by the rotation rate \cite{arri_eb1218_manual}. Magnetic ballasts are
excluded. For an LED fixture, request the dimming method and the PWM
frequency from the supplier; no reviewed source gives a safe value, so the
fixture is accepted only after the on-site test A2.7.

\subsection*{L7. Black surroundings, fixture out of frame.
\textnormal{Decision}}

Black velvet or molton, no light falling on walls. This is how DLR and
Stanford build their testbeds \cite{benninghoff2017epos, park2021tron}, and
black velvet was the best background measured with the same Raspberry Pi HQ
camera this project uses \cite{pauly2022spacelab}. A fixture inside the frame
produces flare; Stanford report exactly that when the camera points near their
sun lamp. IEC 60904-9 separately warns that wall reflections invalidate a
measured uniformity classification \cite{iec60904_9_2020}.

\section{The penumbra test}
\label{sec:penumbra}

Hold a straight edge 1~m in front of a white card placed at the tether
position, and measure the width of the diffuse band with a ruler. The geometry
is
\begin{equation}
w = L \tan \alpha ,
\label{eq:penumbra}
\end{equation}
where $L$ is the edge-to-card distance and $\alpha$ the angular size of the
source as seen from the edge. The measured width therefore gives the angular
size directly. The Sun would give 9.3~mm at $L = 1$~m; the registered limit of
35~mm corresponds to about 2\textdegree; an uncollimated lamp typically gives
87~mm or more. The value of this test is that it converts an optical property
into a ruler reading.

\section{Optics}

\subsection*{O1. Aperture wide open at f/1.6. \textnormal{Decision, with
calculation}}

Two independent arguments point the same way. Stopping down costs light as
the square of the f-number: going to f/2.8 triples the illuminance
requirement, from 2320 to about 7100~lx. And diffraction is already
significant wide open, with the Airy disc at 2.15~\textmu m or 1.39~px at
f/1.6, rising to 2.42~px at f/2.8 and 4.85~px at f/5.6, against a tether that
projects to 1.47~px. There is no aperture at which closing down helps.

\subsection*{O2. Focus at the real tether distance, set by eye.
\textnormal{Physics}}

Set while looking at the live image at full resolution, with the tether in
place. See Chapter \ref{chap:setup} for why the barrel scale cannot be used.

\subsection*{O3. Lock and mark both rings. \textnormal{Critical}}

Tighten the grub screws and mark the ring positions with a pen or tape.
Moving focus changes the focal length, and $f_x$ is the focal length in
pixels, so 1\% of drift is 1\% of scale error on every reconstructed distance,
or 10~mm per metre. It is a systematic bias and averaging does not remove it.
This experiment moves the rig by hand, so a loose ring will move. The marks
are checked at the end of the session, which is criterion A4.4.

\section{Calibration}

\subsection*{C1. Same resolution and same sensor mode as the experiment.
\textnormal{Physics}}

Capture the calibration frames at $4056 \times 3040$ with an explicit
\texttt{-{}-mode}. Without that argument libcamera may silently pick a mode
and then scale or crop, which makes the transfer untrustworthy
\cite{raspberrypi_camera_docs}.

\subsection*{C2. Calibrate on the integrated support. \textnormal{Physics}}

The stereo extrinsics describe the rigid transform between two specific
cameras in specific positions. Calibrating on one mount and then moving the
cameras to another invalidates them. The support is already built; the
calibration session happens on it, in the room, with the structure in place.

\subsection*{C3. Fill all four tilt bins. \textnormal{Physics}}

This is the correction for the diagnosed cause of the 23 previous failures.
With the board nearly parallel to the sensor in every view, a longer focal
length at a greater distance produces an identical image and the focal length
is simply not recoverable, however many frames are taken
\cite{zhang2000flexible, linku_calibration_method}. The guided tool asks for a
specific pose, draws its outline in both previews, and only accepts a frame
that is on that pose, so filling the bins is the tool's job rather than the
operator's; what the operator has to do is physically reach the poses it asks
for, including the strongly tilted ones.

\subsection*{C4. Do not touch the optics afterwards. \textnormal{Critical}}

See O3. Calibration and experiment are one continuous optical configuration.

\section{Camera settings during the checks and the run}

Base gain, fixed white balance, manual exposure at the ceiling given in Table
\ref{tab:optics-at-work}. Automatic exposure and automatic white balance are
both disabled: an auto-exposure loop will chase the moving tether and change
the radiometry from frame to frame, which breaks both the flicker test and the
stereo consistency the matching depends on.
